\section{Discussion}

Previous work showed that R2 and R4 neurons are required for the lifespan-shortening effects of death perception, establishing these cells as a neural substrate through which sensory experience can influence aging \citep{gendron2023}. 
Here, we show that direct activation of R2 neurons, co-activation of R2/R4 neurons, and inhibition of R4d neurons are each sufficient to promote longevity under standard laboratory conditions in adult flies of both sexes. 
Strikingly, these effects are elicited by manipulating very small neuronal populations---in some cases, fewer than ten neurons per hemisphere \citep{omoto2018}---underscoring the sensitivity of organismal aging to defined circuit states. 
These findings argue that the EB acts instructively rather than merely as a passive conduit for lifespan-altering sensory information.

We also show that distinct EB manipulations converge on increased lifespan through divergent molecular states, as bulk head transcriptomics after one week of treatment revealed that R2 activation, R4d inhibition, and R2/R4 co-activation did not produce a common downstream transcriptional program. 
R4d inhibition produced a more limited transcriptional response than R2 or R2/R4 activation, but attenuated stress-associated changes observed in green-light controls. 
Temperatures in our ventilated optogenetic lifespan apparatuses were not significantly higher than those experienced by flies aged in constant darkness, indicating that something other than heat was the primary cause of the effects of green light on lifespan and gene expression. 
Fly vision is more sensitive to green than to red light \citep{salcedo1999}, and the short lifespan (\figref{fig:lifespanfig}C) and over-representation of upregulated heat shock proteins suggest that it represents a meaningful stressor at the intensity used in our experiments to ensure deep brain penetration. 
If so, it may be that a muted molecular response to moderate long-term stress contributed to the extended lifespan of flies with silenced R4d neurons. 
It is, however, currently unknown whether similar mechanisms influence lifespan in more benign conditions. 

The transcriptional data also point to specific pathways that may link EB activity to aging-regulatory physiology. 
The induction of \textit{foxo}, \textit{Thor}, \textit{InR}, \textit{dilp6}, and \textit{Cbs} following R2 activation suggests engagement of nutrient-sensing, proteostatic, and sulfur-amino-acid metabolic processes previously implicated in lifespan control. 
Notably, however, the R2-induced state does not resemble a canonical starvation program. 
Nutrient-responsive IIS components were induced alongside ribosome biogenesis and translational programs that are not ordinarily co-regulated in this way \citep{essers2016}, suggesting that direct EB activation may uncouple pathways that are typically coordinated by nutrient availability. 
Because prior work showed that R2/R4 neurons regulate \textit{dilp3} and \textit{dilp5} expression in insulin-producing cells during death perception \citep{gendron2023}, one plausible hypothesis is that EB-dependent internal states are translated into systemic aging effects through downstream neuroendocrine centers. 
However, the present data do not establish that mechanism directly, and defining the relevant effector pathways will require future work.

Notably, the molecular signatures induced by EB manipulation did not manifest as simple changes in behavior or physiology. 
Our GO analysis revealed markedly different transcriptional indicators of nutrient availability in the brain following R2 and R2/R4d activation. 
General starvation-response processes were upregulated in R2-activated heads (\figref{fig:seqfig}B).
One might predict such changes to result from actual nutrient scarcity if, for example, flies were consuming less. 
If, on the other hand, these transcriptional signatures reflected a hunger-like neural state, one might expect to observe increased feeding to satisfy that perceived demand. 
We observed neither, arguing that the transcriptomic changes are unlikely to arise secondarily from altered food consumption.
By contrast, starvation-response processes were downregulated in R2/R4d-activated heads (\figref{fig:seqfig2}C). 
Activation instead produced a transcriptome enriched for biosynthetic and storage-associated processes, including protein synthesis and glycogen metabolism, suggestive of a nutrient-replete state. 
However, we again saw no corresponding change in feeding that would be expected if this state were behaviorally driven. 
Together, these findings indicate that the transcriptional signatures elicited by these manipulations are a direct result of neural activation itself rather than a secondary effect of changes in behavior. 
More broadly, they argue against dietary restriction as the basis of lifespan extension in either manipulation and suggest that EB circuit activation can impose internal nutrient-state programs while uncoupling those states from the motivational and behavioral outputs that normally regulate feeding. 

Despite unchanged feeding, R2 males and both R2/R4 sexes had reduced TAG stores. 
These TAG changes did not consistently track lifespan extension or starvation resistance, and were not caused by increased locomotor activity. 
Negative geotaxis similarly failed to reveal early behavioral correlates of the transcriptional states measured after one week of manipulation. 
Together, these results indicate that EB-dependent regulation of aging is not readily explained by gross changes in feeding, energy expenditure, lipid depletion, or locomotor performance.

Although EB manipulations imposed distinct molecular states without corresponding differences in gross behavior or physiology, valence strongly distinguished them from one another. 
In the FLIC, R2 activation in females was aversive, promoting rapid disengagement from the sucrose well without suppressing re-approach, while R2/R4 co-activation was positively valenced, increasing interactions and return frequency without prolonging individual feeding bouts. 
In a feeding-independent place-preference assay, the positive valence of R2/R4 co-activation generalized across sex and stimulation regime. 
Interestingly, the negative valence of R2 and R4d activation was more context dependent, being restricted to females across both light regimes (R2), or restricted to one light regime but generalizing across sex (R4d). 
These findings indicate that EB ring neurons do not contribute fixed, additive aversive or appetitive signals. 
Rather, the motivational meaning of EB activity appears to depend on population combination, assay context, sex, and stimulation conditions. 
Particularly striking is the observation that both aversive and appetitive EB states can promote longevity. 
This argues against a simple model in which lifespan extension is coupled to one privileged valence state and instead suggests that distinct internal states may engage different downstream mechanisms that converge on slower aging.

This nonlinear relationship between EB population activity and behavioral output may be especially informative in light of known interactions among ring-neuron classes. 
Prior work suggests that R2, R4m, and R4d neurons participate in an incoherent feedforward circuit, raising the possibility that manipulations such as R2 activation, R2/R4 co-activation, and R4d inhibition may overlap in some network-level consequences even while producing distinct transcriptional and behavioral states. 
Our data do not resolve whether these interventions ultimately converge on a shared terminal mechanism or instead reach similar lifespan outcomes through partially distinct pathways. 
However, the marked divergence in both transcriptional state and valence across manipulations that each extend lifespan suggests that the relationship between EB activity and aging is not fixed to a single circuit configuration or downstream program.
A more parsimonious interpretation is that EB circuits influence aging through population-level configurations whose downstream consequences depend on the particular combination of neurons engaged.

Profiling of peripheral tissues revealed that R2 activation drove coordinated systemic remodeling beyond the brain. 
In bodies, R2 activation suppressed proliferative and DNA repair programs while upregulating a vacuolar H$^{+}$-ATPase-dominated electrochemical homeostasis program not observed in heads, indicating that EB activity orchestrates tissue-specific rather than uniform physiological responses. 
This V-ATPase program was not accompanied by the upregulation of autophagy or trafficking pathways, suggesting a decoupling of organellar acidification from degradative or secretory functions. 
The upregulation of inwardly rectifying potassium channels may serve to maintain cell membrane potential during increased organellar proton flux, and the enrichment of fatty acid and ether phospholipid synthesis genes suggests a remodeling of cellular and organellar membranes to support increased ion transport. 
Metabolomic profiling reinforced this interpretation: the depletion of free nucleotides and amino acids likely reflects the reduced proliferative burden observed in our transcriptional data, while reductions in choline, glycerophosphocholine, and acylcarnitines may reflect the increase in membrane remodeling activity. 
Among upregulated metabolites, the accumulation of glycolytic and TCA-cycle intermediates may reflect reduced demand for proliferative inputs, while glutamine and arginine may likewise accumulate as a consequence of reduced demand from growth-related processes. 
Serine and betaine serve as methyl donors, and their elevation, along with depletions in SAH and choline and the upregulation of both \textit{Sam-S} and \textit{Gnmt} in head and body, indicates a restructuring of methionine cycle flux. 
Although none of the core kynurenine pathway genes were upregulated in our RNA-seq dataset, the coordinated elevation of kynurenic acid, quinolinic acid, and anthranilic acid could account for the depleted tryptophan signature. 
The convergence of transcriptomic and metabolomic datasets indicates an orchestrated remodeling of peripheral physiology imposed by the activation of a small number of R2 neurons in the brain, one that does not resemble canonical starvation or dietary restriction and further supports the conclusion that EB neuron activation imposes a distinct physiological state rather than recapitulating a known nutritional condition.

Our findings that discrete or combined populations of ellipsoid body ring neurons can instructively modulate lifespan through the generation of distinct internal states carry potentially significant implications for understanding the neural basis of aging in humans. 
The \textit{Drosophila} central complex, of which the ellipsoid body is a core component, shares strong neuroanatomical and functional homology with the vertebrate basal ganglia \citep{strausfeld2013}. 
Both structures serve as integrative hubs that process sensory information, coordinate motor output, and critically contribute to the generation of motivational and affective states. 
In particular, limbic basal ganglia output pathways, including ventral pallidal circuits, can assign appetitive or aversive value to ongoing behavior and reshape responses to challenge in a state-dependent manner \citep{stephenson-jones2020}. 
Our findings raise the possibility that a similarly general principle applies to aging: circuit configurations that impose distinct motivational or physiological states may converge on long-term effects on organismal maintenance and survival. 
More generally, these results suggest that lifespan may be modulated not only by what animals perceive, but by how neural state transforms the meaning of those perceptions into coordinated physiological output. 
If analogous circuits in the human brain similarly encode valence states that influence systemic physiology, this would provide a mechanistic framework for understanding associations between psychological states, chronic stress, and aging in humans.

\subsubsection*{Limitations}
Several limitations should be considered when interpreting these findings. 
First, our transcriptomic analyses were performed on female heads at a single early time point, and thus do not capture male-specific or later-emerging molecular states. 
Second, bulk head RNA-seq cannot distinguish transcriptional changes arising within the brain from those occurring in peripheral head tissues, nor can it resolve cell-type-specific responses downstream of EB manipulation. 
Third, the green light required for GtACR1 complicated behavioral interpretation of some inhibition experiments, particularly in the FLIC, where abrupt light onset itself appeared to be startling. 
Newly developed red-responsive anion channelrhodopsins will help mitigate this in future experiments \citep{bushey2025}. 
Finally, although the present work shows that EB manipulations are sufficient to alter lifespan and to impose distinct transcriptional and motivational states, it does not yet establish the causal chain linking EB activity to the peripheral transcriptional and metabolic changes observed, nor identify the downstream endocrine effectors that execute these aging phenotypes. 
These limitations define clear priorities for future work, including cell-type-resolved transcriptional analysis, direct study of neural insulin-producing cells and neuroendocrine outputs, and more refined dissection of sex-specific responses.